% Generated draft: requires mathematical and source review.
\documentclass[11pt]{article}
\usepackage[margin=25mm]{geometry}
\usepackage{fontspec}
\setmainfont{Latin Modern Roman}
\usepackage{amsmath,amssymb,mathtools,unicode-math}
\setmathfont{Latin Modern Math}
\usepackage{enumitem,needspace,xurl,hyperref}
\setlength{\parindent}{0pt}
\setlength{\parskip}{5pt}
\setlength{\emergencystretch}{3em}
\newcommand{\N}{\mathbb{N}}
\newcommand{\Z}{\mathbb{Z}}
\newcommand{\Q}{\mathbb{Q}}
\newcommand{\R}{\mathbb{R}}
\newcommand{\C}{\mathbb{C}}
\newcommand{\F}{\mathbb{F}}
\newcommand{\A}{\mathbb{A}}
\newcommand{\PP}{\mathbb{P}}
\newcommand{\E}{\mathbb{E}}
\newcommand{\eps}{\varepsilon}
\newcommand{\dd}{\,\mathrm{d}}
\begin{document}
\section*{UnsolvedMath ID: 20000463}
\textbf{Source problem code:} catalogue entry
\subsection{Definitions and mathematical statement}
The complete current mathematical-source entry is: ``Question 30. This problem has been withdrawn.'' No active mathematical question, assumptions, or target statement is supplied. The trailing extracted ``Problem/'' is layout context. This withdrawal record is preserved under its permanent identity rather than replaced by a conjecture inferred from the attached research aid.

For historical context only, the report says an official pre-withdrawal archive named two cubic surfaces. It identifies the first exactly with Question 5's \(D_4\) surface and the second with Question 5(vii)'s \(E_6\) surface after \(z\mapsto-z\). It additionally gives the projective transformation
\[
[x:y:z:t]\longmapsto[y:x:t:-z]
\]
as a \(GL_4(\Z)\) isomorphism from the latter surface to a published standard \(E_6\) model. As a signed permutation of primitive integral coordinates, this preserves maximum height \(H([x:y:z:t])=\max(|x|,|y|,|z|,|t|)\). The actual surface equations and archive identifier are not included, so the equation equivalence cannot be checked from this record.

Under the report's maximum-height, line-deleted interpretation of the historical point-count questions, it says Le Boudec's 2014 theorem and the 2007 theorem of de la Bret\`eche--Browning--Derenthal settle the respective problems. No theorem statements or full references are supplied. Exact redundancy is offered as a plausible explanation for withdrawal, but the report explicitly says the editorial motive is undocumented. None of this historical context reactivates the withdrawn question.
\subsection{Short English statement}
This is a withdrawn-question record, not an active mathematical problem. The attached report offers historical duplicate identifications and a height-preserving coordinate map, without the original equations or a documented withdrawal motive.
\subsection{Sources}
\begin{itemize}
\item AIM workshop ``Rational and integral points on higher dimensional varieties,'' Question 30: \url{https://aimath.org/WWN/qptsurface2/qptsurface2.pdf}. The supplied current text says the problem has been withdrawn; no PDF was accessed here.
\item Local canonical registry \texttt{aim-algebraic-number-theory-notes.json}, \texttt{notes[114]}, tag \texttt{question}; permanent UnsolvedMath ID \(20000463\), source code \texttt{catalogue entry}. Original remarks and literature field are empty.
\item Supplied research attempt 1 refers to an unidentified official pre-withdrawal archive, Le Boudec's 2014 result, and de la Bret\`eche--Browning--Derenthal's 2007 result, without full authoritative citations or surface equations. It is a machine-generated, non-peer-reviewed research aid requiring independent verification before citation.
\end{itemize}
\end{document}
